% Figure 13-B -- full CCDL schematic of CCB-D in drawn component symbols.
% CJTs use the Chapter 3 symbol (NPN in a circle) rotated 90 deg CCW, so that
% cause enters from the left, bias enters from below, and effect leaves to the
% right -- the orientation that suits a left-to-right signal path over a bottom
% bias rail. Q0's three terminals are named on the drawing.
\begin{tikzpicture}[
  x=1cm, y=1cm,
  lbl/.style={font=\sffamily\scriptsize, align=center, inner sep=1pt},
  llbl/.style={font=\sffamily\scriptsize, align=left, inner sep=1pt},
  tlbl/.style={font=\sffamily\scriptsize\itshape, align=center, inner sep=1pt},
  loopbox/.style={nblock, font=\sffamily\scriptsize, minimum height=6mm,
                  inner sep=2.5pt},
  wire/.style={line width=0.6pt, draw=black},
  % ---- component symbols -------------------------------------------------
  cjt/.pic={
    \draw[line width=1.0pt] (-0.22,-0.10) -- (0.22,-0.10);
    \draw[line width=0.6pt] (0,-0.10) -- (0,-0.55);
    \draw[line width=0.6pt] (-0.10,-0.10) -- (-0.30,0.20) -- (-0.60,0.20);
    \draw[-{Latex[length=1.5mm,width=1.2mm]}, line width=0.6pt]
      (0.10,-0.10) -- (0.245,0.1175);
    \draw[line width=0.6pt] (0.245,0.1175) -- (0.30,0.20) -- (0.60,0.20);
    \draw[line width=0.6pt] (0,0.02) circle (0.40);
  },
  diode/.pic={
    \filldraw[fill=white, draw=black, line width=0.6pt]
      (-0.17,-0.18) rectangle (0.17,0.18);
    \filldraw[fill=white, draw=black, line width=0.6pt]
      (-0.17,-0.18) -- (-0.17,0.18) -- (0.11,0) -- cycle;
    \draw[line width=0.9pt] (0.11,-0.20) -- (0.11,0.20);
  },
  diodeL/.pic={
    \filldraw[fill=white, draw=black, line width=0.6pt]
      (-0.17,-0.18) rectangle (0.17,0.18);
    \filldraw[fill=white, draw=black, line width=0.6pt]
      (0.17,-0.18) -- (0.17,0.18) -- (-0.11,0) -- cycle;
    \draw[line width=0.9pt] (-0.11,-0.20) -- (-0.11,0.20);
  },
  fuse/.pic={
    \draw[line width=0.6pt] (0,-0.30) -- (0,0.30);
    \filldraw[fill=white, draw=black, line width=0.6pt]
      (-0.13,-0.18) rectangle (0.13,0.18);
    \draw[line width=0.6pt] (0,-0.18) -- (0,0.18);
  },
  res/.pic={
    \filldraw[fill=white, draw=black, line width=0.6pt]
      (-0.13,-0.25) rectangle (0.13,0.25);
  },
  gnd/.pic={
    \draw[line width=0.7pt] (-0.20,0) -- (0.20,0);
    \draw[line width=0.7pt] (-0.13,-0.10) -- (0.13,-0.10);
    \draw[line width=0.7pt] (-0.06,-0.20) -- (0.06,-0.20);
  },
  % capacitor with its return to ground; origin at the top terminal
  capgnd/.pic={
    \draw[line width=0.6pt] (0,0) -- (0,-0.35);
    \draw[line width=0.9pt] (-0.23,-0.35) -- (0.23,-0.35);
    \draw[line width=0.9pt] (-0.23,-0.52) -- (0.23,-0.52);
    \draw[line width=0.6pt] (0,-0.52) -- (0,-0.82);
    \draw[line width=0.7pt] (-0.20,-0.82) -- (0.20,-0.82);
    \draw[line width=0.7pt] (-0.13,-0.92) -- (0.13,-0.92);
    \draw[line width=0.7pt] (-0.06,-1.02) -- (0.06,-1.02);
  }
]

% ================= STAGE 1: input, D1, R1 divider =======================
\node[casterm, anchor=west] at (0,2.75) {\casnote{\detokenize{^0_desire^C_{F2}|A}}};
\draw[wire] (0.95,2.45) -- (0.95,1.90);
\fill (0.95,1.90) circle (0.05);

\node[casterm, anchor=west] at (2.60,1.90) {\casnote{\detokenize{^P-origin}}};
\draw[wire] (2.60,1.90) -- (0.95,1.90);
\pic at (2.15,1.90) {diodeL};
\node[lbl, anchor=south] at (2.15,2.10) {\textbf{D1}\ \ Vf 0.7};

\draw[wire] (0.95,1.90) -- (0.95,1.55);
\pic at (0.95,1.30) {res};
\node[llbl, anchor=west] at (1.18,1.30) {\textbf{R1a}\ \ \casnote{\detokenize{Dial(^2)}}};
\draw[wire] (0.95,1.05) -- (0.95,-0.45);
\fill (0.95,0) circle (0.05);
\node[lbl, anchor=south] at (1.38,0.06) {tap $k$};
\pic at (0.95,-0.70) {res};
\draw[wire] (0.95,-0.95) -- (0.95,-1.25);
\pic at (0.95,-1.25) {gnd};
\node[lbl, anchor=north] at (0.95,-1.58)
  {\textbf{R1b}\ \ \casnote{\detokenize{Dial(^3)}}};

% ================= STAGE 2: MUX1 in the main path ======================
\draw[wire] (0.95,0) -- (2.32,0) -- (2.32,0.24) -- (2.50,0.24);
\fill (1.80,0) circle (0.05);

\filldraw[fill=casCfill, draw=casCedge, line width=0.6pt]
  (2.50,0.50) -- (3.35,0.26) -- (3.35,-0.26) -- (2.50,-0.50) -- cycle;
\node[font=\sffamily\scriptsize] at (2.92,0) {MUX1};
\node[llbl, anchor=north west] at (3.44,-0.30) {pass if\\ TP2 $<$ 0.30};

% ================= STAGE 3: D2 and Q2 (AMPLIFY) ========================
\draw[wire] (3.35,0) -- (4.95,0);
\pic at (4.10,0) {diode};
\node[lbl, anchor=south] at (4.10,0.26)
  {\textbf{D2}\ \ Vf 1.2\\ \casnote{\detokenize{^C -> ^E}}};

\pic at (5.55,-0.20) {cjt};
\draw[wire] (5.55,-0.75) -- (5.55,-3.90);
\node[lbl, anchor=south] at (5.55,0.86) {\textbf{Q2}\ \ $\beta$ = by type};
\draw[line width=0.7pt, dashed] (5.30,0.46) -- (6.60,0.46);
\node[lbl, anchor=south] at (5.95,0.50) {ceiling\ \ sat = 0.80};

% ================= STAGE 4: OSC1 (RHYTHM) ==============================
\draw[wire] (6.15,0) -- (6.60,0);
\filldraw[fill=casEfill, draw=casEedge, line width=0.6pt]
  (6.60,-0.45) rectangle (7.70,0.45);
\draw[line width=0.7pt]
  (6.76,0.10) -- (6.76,0.32) -- (6.95,0.32) -- (6.95,0.10) --
  (7.15,0.10) -- (7.15,0.32) -- (7.34,0.32) -- (7.34,0.10) -- (7.54,0.10);
\node[font=\sffamily\scriptsize] at (7.15,-0.22) {OSC1};
\node[llbl, anchor=north west] at (7.32,-0.52)
  {24\,hr period\\ duty $\le$ 0.05};

% ================= STAGE 5: D3 and Q3 (MANIFEST) =======================
\draw[wire] (7.70,0) -- (8.85,0);
\pic at (8.20,0) {diode};
\node[lbl, anchor=south] at (8.20,0.52)
  {\textbf{D3}\ \ Vf 2.0\\ \casnote{\detokenize{^E -> ^O}}};

\pic at (9.45,-0.20) {cjt};
\node[lbl, anchor=south west] at (9.05,0.26)
  {\textbf{Q3}\\ $\beta$ = recep.\\ \casnote{\detokenize{^9^O}}};
\draw[wire] (10.05,0) -- (10.30,0);
\fill (10.15,0) circle (0.05);
\filldraw[fill=white, draw=black, line width=0.6pt] (10.30,0) circle (0.07);

% Q3's bias is NOT the operator rail -- it is a separate collective store.
\draw[wire] (9.45,-0.75) -- (9.45,-1.45);
\pic at (9.45,-1.45) {capgnd};
\node[lbl, anchor=east] at (9.35,-1.45) {\casnote{\detokenize{^5_collective^O}}};

% ================= POLICY: audit chain Q0 -> Q1 -> F1 ==================
\draw[wire] (1.80,0) -- (1.80,-2.35) -- (2.00,-2.35);
\pic at (2.60,-2.55) {cjt};
\node[lbl, anchor=south] at (2.60,-2.10) {\textbf{Q0}\\ $\beta$ = 1.0};
\draw[wire] (2.60,-3.10) -- (2.60,-3.90);

\draw[wire] (3.20,-2.35) -- (3.60,-2.35);
\pic at (4.20,-2.55) {cjt};
\node[lbl, anchor=south] at (4.20,-2.10)
  {\textbf{Q1}\ \ $\beta$ = 2.0\\ sat 0.85\ \ cut 0.30};
\draw[wire] (4.20,-3.10) -- (4.20,-3.90);

\draw[wire] (4.80,-2.35) -- (5.15,-2.35);
\fill (5.15,-2.35) circle (0.05);
\draw[wire] (5.15,-2.35) -- (5.15,-3.20);
\pic at (5.15,-2.82) {fuse};
\node[lbl, anchor=west] at (5.34,-2.82) {\textbf{F1}\\ 0.90};
\pic at (5.15,-3.20) {gnd};

% terminal names, taught once on Q0
\node[tlbl, anchor=north east] at (2.15,-2.43) {cause};
\node[tlbl, anchor=north]      at (3.40,-2.45) {effect};
\node[tlbl, anchor=west]       at (2.70,-3.01) {bias};

% refusal path back to MUX1: second input, and the select line
\draw[wire] (5.15,-1.05) -- (2.12,-1.05) -- (2.12,-0.24) -- (2.50,-0.24);
\fill (2.90,-1.05) circle (0.05);
\draw[flow] (2.90,-1.05) -- (2.90,-0.39);
\node[lbl, anchor=east] at (2.82,-0.72) {sel};

% ================= bias rail and C1 ====================================
\draw[line width=1.3pt] (0.55,-3.90) -- (5.55,-3.90);
\foreach \x in {2.60,4.20,5.55}{\fill (\x,-3.90) circle (0.05);}
\node[lbl, anchor=west] at (6.15,-3.90) {\casnote{\detokenize{^5^P}} bias rail};

\pic at (0.55,-3.90) {capgnd};
\node[llbl, anchor=west] at (0.95,-4.30)
  {\textbf{C1}\ \ \casnote{\detokenize{^5_operator^P_{F7}|T}}\\
   $\tau$ = 3333\,hr,\ leak 0.0002/day};

% ================= LOOP A: FLT1, fast, stays inside ^O =================
\draw[fbflow, rounded corners=2pt]
  (10.15,0) -- (10.15,-2.95) -- (6.70,-2.95) -- (6.70,-0.45);
\node[loopbox] at (8.55,-2.95) {FLT1\quad $\tau$ = 24\,hr};

% ================= LOOP B: FLT2 + FB1, slow, ends at C1 ================
\draw[fbflow, rounded corners=2pt] (10.15,-2.95) -- (10.15,-5.40) -- (8.85,-5.40);
\draw[fbflow] (6.75,-5.40) -- (5.20,-5.40);
\draw[fbflow, rounded corners=2pt]
  (3.40,-5.40) -- (0.15,-5.40) -- (0.15,-3.90) -- (0.50,-3.90);
\node[loopbox] at (7.80,-5.40) {FLT2\quad $\tau$ = 3\,weeks};
\node[loopbox] at (4.30,-5.40) {FB1\quad $\beta$ = 0.05};
\node[tlbl, anchor=north] at (4.30,-5.72)
  {terminates at the bias terminal, never at the signal path};

\end{tikzpicture}
